\chapter{Sistema de UI/UX}

\section{Tecnolog\'ias}

\begin{itemize}[leftmargin=*]
    \item \textbf{TailwindCSS v4:} Framework de estilos utility-first con
          configuraci\'on via CSS.
    \item \textbf{shadcn/ui:} Componentes base accesibles construidos sobre
          Radix UI.
    \item \textbf{Framer Motion:} Animaciones declarativas con API de spring.
    \item \textbf{next-themes:} Soporte de tema claro/oscuro/sistema.
    \item \textbf{lucide-react:} Iconos SVG consistentes.
    \item \textbf{Radix UI:} Primitivos accesibles para UI.
\end{itemize}

\section{Sistema de Temas de Color}

El archivo \texttt{globals.css} define 6 temas de color mediante variables
CSS personalizadas:

\begin{lstlisting}[style=css, caption=Temas de color en CSS]
[data-theme-color="blue"] {
  --primary: 221.2 83.2% 53.3%;
  --primary-foreground: 210 40% 98%;
  --ring: 221.2 83.2% 53.3%;
}

[data-theme-color="terracota"] {
  --primary: 12 76% 61%;
  --primary-foreground: 0 0% 100%;
  --ring: 12 76% 61%;
}

[data-theme-color="pink"] {
  --primary: 336 80% 58%;
  --primary-foreground: 0 0% 100%;
  --ring: 336 80% 58%;
}

[data-theme-color="purple"] {
  --primary: 271 81% 56%;
  --primary-foreground: 0 0% 100%;
  --ring: 271 81% 56%;
}

[data-theme-color="green"] {
  --primary: 142 76% 36%;
  --primary-foreground: 0 0% 100%;
  --ring: 142 76% 36%;
}

[data-theme-color="orange"] {
  --primary: 24 95% 53%;
  --primary-foreground: 0 0% 100%;
  --ring: 24 95% 53%;
}
\end{lstlisting}

\subsection{ThemeColorSwitcher}

El componente \texttt{ThemeColorSwitcher} es un bot\'on flotante (FAB) en
la esquina inferior derecha que:

\begin{itemize}[leftmargin=*]
    \item Muestra 6 opciones de color con indicador de selecci\'on activa.
    \item Usa Framer Motion para animaciones de apertura/cierre.
    \item Almacena la preferencia mediante un custom hook
          (\texttt{use-theme-color}).
    \item Cambia el atributo \texttt{data-theme-color} en el documento.
\end{itemize}

\section{Tema Claro/Oscuro}

\begin{itemize}[leftmargin=*]
    \item Implementado con \texttt{next-themes}.
    \item Soporta tema claro, oscuro y seguimiento del sistema.
    \item Bot\'on de toggle en la Navbar con iconos Sun/Moon.
    \item Variables CSS en \texttt{:root} (claro) y \texttt{.dark} (oscuro).
\end{itemize}

\section{Componentes shadcn/ui}

Los componentes UI base siguen la convenci\'on \texttt{cn()} para fusi\'on
de clases:

\begin{lstlisting}[style=typescript, caption=Funcion cn() para fusion de clases]
import { type ClassValue, clsx } from 'clsx';
import { twMerge } from 'tailwind-merge';

export function cn(...inputs: ClassValue[]) {
  return twMerge(clsx(inputs));
}
\end{lstlisting}

Componentes implementados:

\begin{longtable}{|l|l|}
\hline
\textbf{Componente} & \textbf{Caracter\'isticas} \\
\hline
Button & 6 variantes (default, destructive, outline, secondary,
         ghost, link), 4 tama\~nos (default, sm, lg, icon) \\
Card & Header, content, footer con sem\'antica \\
Input & Campo de texto con validaci\'on \\
Label & Etiqueta accesible para formularios \\
Avatar & Imagen circular con fallback de iniciales \\
Separator & L\'{\i}nea divisoria horizontal \\
DropdownMenu & Men\'u contextual desplegable con items \\
Textarea & \'Area de texto multil\'{\i}nea \\
\hline
\end{longtable}

\section{Layout Responsivo}

\subsection{Sidebar}

\begin{itemize}[leftmargin=*]
    \item Colapsable en dispositivos m\'oviles (estado local con toggle).
    \item Animaciones con Framer Motion y \texttt{layoutId} para
          transiciones fluidas.
    \item Iconos de mascota decorativos semi-transparentes.
    \item Navegaci\'on con resaltado de ruta activa.
    \item Contenido principal con sidebar, navbar y \'area de contenido.
\end{itemize}

\subsection{Navbar}

\begin{itemize}[leftmargin=*]
    \item Campana de notificaciones con polling cada 15 segundos y
          badge de contador.
    \item Toggle de tema claro/oscuro.
    \item Men\'u de usuario con avatar, iniciales y cierre de sesi\'on.
\end{itemize}

\section{Animaciones}

\subsection{Animaciones Personalizadas (CSS)}

\begin{lstlisting}[style=css, caption=Animaciones CSS personalizadas]
@keyframes float {
  0%, 100% { transform: translateY(0px); }
  50% { transform: translateY(-10px); }
}

@keyframes paw-print {
  0% { opacity: 0; transform: scale(0) rotate(0deg); }
  50% { opacity: 0.5; transform: scale(1.2) rotate(10deg); }
  100% { opacity: 0; transform: scale(1) rotate(0deg); }
}

@keyframes gradient-shift {
  0% { background-position: 0% 50%; }
  50% { background-position: 100% 50%; }
  100% { background-position: 0% 50%; }
}
\end{lstlisting}

\subsection{Framer Motion}

Las animaciones con Framer Motion se utilizan en:

\begin{itemize}[leftmargin=*]
    \item Transiciones de p\'agina con \texttt{staggerChildren}.
    \item Sidebar con \texttt{layout} y \texttt{layoutId} para animaciones
          de colapso/expansi\'on.
    \item Tarjetas de dashboard con hover (scale + y-offset).
    \item P\'agina de login con mascotas flotantes, gradientes animados
          y entradas escalonadas.
    \item Indicador activo de sidebar con animaci\'on spring.
\end{itemize}

\section{Layout Principal}

El layout del dashboard usa route groups de Next.js:

\begin{lstlisting}[style=typescript, caption=DashboardLayout - Sidebar + Navbar]
export default function DashboardLayout({
  children
}: { children: React.ReactNode }) {
  return (
    <div className="flex min-h-screen">
      <Sidebar />
      <div className="flex flex-1 flex-col">
        <Navbar />
        <main className="flex-1 p-6">{children}</main>
      </div>
      <ThemeColorSwitcher />
    </div>
  );
}
\end{lstlisting}

\section{P\'agina de Login}

La p\'agina de login presenta un dise\~no de dos columnas:

\begin{itemize}[leftmargin=*]
    \item \textbf{Lado izquierdo (solo escritorio):} Marca con iconos de
          mascota flotantes, nombre del sistema, tarjetas de caracter\'isticas
          y huellas animadas. Fondo con gradiente animado.
    \item \textbf{Lado derecho:} Formulario de inicio de sesi\'on con
          animaciones de entrada escalonadas, selector de visibilidad de
          contrase\~na, credenciales de prueba y bot\'on con spinner de carga.
\end{itemize}
